\chapter{Experimental Results}

This chapter presents the results of the 4,080-run benchmark comparing the full simulation engine (\engine) against the naive baseline (\naive) across 20 real-world scenarios, three actor scales, and four repetitions per configuration. All reported confidence intervals are at the 95\% level, and all statistical tests employ a paired design unless otherwise noted.


\section{Overall Performance Comparison}

Table~\ref{tab:overall} presents the aggregate results across all 4,080 simulation runs. The engine achieves lower persona drift, fewer envelope violations, higher commitment fulfillment, and greater dialogue coherence than the naive baseline across all primary and secondary metrics.

\begin{table}[t]
\centering
\caption{Overall Performance Comparison ($N = 4{,}080$ runs)}
\label{tab:overall}
\begin{tabular}{lcc}
\toprule
Metric & \engine & \naive \\
\midrule
\multicolumn{3}{l}{\textit{Primary Metrics}} \\
Persona Drift MAE       & \textbf{0.1678 $\pm$ 0.017} & 0.1780 $\pm$ 0.018 \\
Envelope Violations     & \textbf{2.057}              & 2.283 \\
95\% CI (Drift)         & [0.167, 0.169]              & [0.177, 0.179] \\
\midrule
\multicolumn{3}{l}{\textit{Commitment \& Consistency}} \\
Contradiction Rate      & 0.000                       & 0.000 \\
Fulfillment Rate        & \textbf{0.781}              & 0.675 \\
Semantic Identity       & \textbf{0.411}              & 0.402 \\
\midrule
\multicolumn{3}{l}{\textit{Dialogue Quality}} \\
Dialogue Coherence      & \textbf{0.106}              & 0.083 \\
Topic Drift Rate        & \textbf{0.521}              & 0.541 \\
Repetition Rate         & 0.002                       & 0.000 \\
\midrule
\multicolumn{3}{l}{\textit{Action Layer (engine only)}} \\
Action Validity         & 0.641                       & --- \\
Action-Plan Alignment   & 0.968                       & --- \\
Planned Action Coverage & 0.928                       & --- \\
Role Action Diversity   & 0.510                       & --- \\
\bottomrule
\end{tabular}
\end{table}

The engine reduces persona drift by 5.7\% relative to the naive baseline (0.1678 vs.\ 0.1780). While this improvement may appear modest in absolute terms, the 95\% confidence intervals are non-overlapping: the engine's drift CI of $[0.167, 0.169]$ does not intersect with the naive CI of $[0.177, 0.179]$, confirming that the difference is statistically significant. The paired $t$-test yields $t = -10.24$, $p < 0.0001$, with Cohen's $d = -0.94$ (large effect).

Envelope violations are reduced by 9.9\% (2.057 vs.\ 2.283), indicating that the engine not only reduces average trait deviation but also decreases the frequency of significant personality boundary breaches.

Commitment contradiction is 0.000 under both conditions, reflecting the effectiveness of the Commitment Lifecycle component. However, commitment fulfillment---the proportion of stated commitments that are subsequently acted upon---is substantially higher under the engine (0.781 vs.\ 0.675), representing a 15.7\% improvement. This suggests that the engine's soft prompt injection (``Your prior positions:'') not only prevents contradictions but also promotes follow-through.

Dialogue coherence, measured as the semantic similarity between consecutive turns by the same actor, is 27.8\% higher under the engine (0.106 vs.\ 0.083), indicating that the candidate selection process favors responses that maintain topical continuity.


\section{Per-Trait OCEAN Analysis}

Table~\ref{tab:pertrait} disaggregates persona drift by individual OCEAN trait, revealing that the engine's improvement is concentrated in specific traits while others remain resistant to intervention.

\begin{table}[t]
\centering
\caption{Per-Trait Absolute Error by Condition}
\label{tab:pertrait}
\begin{tabular}{lccc}
\toprule
Trait & \engine & \naive & $\Delta$ \\
\midrule
Openness (O)          & \textbf{0.165} & 0.183 & $-$10.0\% \\
Conscientiousness (C) & 0.221          & 0.224 & $-$1.6\% \\
Extraversion (E)      & \textbf{0.120} & 0.144 & $-$16.5\% \\
Agreeableness (A)     & \textbf{0.163} & 0.171 & $-$4.6\% \\
Neuroticism (N)       & 0.170          & 0.168 & +1.2\% \\
\bottomrule
\end{tabular}
\end{table}

\textbf{Openness and Extraversion} exhibit the largest improvements ($-10.0\%$ and $-16.5\%$ respectively). These traits have the strongest psycholinguistic signals in the 30-feature extraction pipeline: Extraversion correlates with word count, question frequency, and turn initiation patterns, while Openness correlates with idea generation and hypothetical reasoning. The engine's candidate selection process can effectively discriminate among candidates on these well-signaled traits.

\textbf{Conscientiousness} shows minimal improvement ($-1.6\%$, from 0.224 to 0.221). This trait has the highest absolute error of all five traits under both conditions, reflecting a persistent estimation difficulty. The dynamic calibration introduced in earlier experiments reduced C-error by 38\% compared to static coefficients, but the remaining error appears to represent a measurement floor: Conscientiousness markers (planning language, structure markers, action items) are frequently used by all actors in policy deliberation regardless of their assigned C level, making differential estimation difficult.

\textbf{Neuroticism} shows no improvement and even a slight regression (+1.2\%, from 0.168 to 0.170). This result is the most theoretically significant finding of the per-trait analysis. The \texttt{self\_doubt\_count} feature---the most direct marker of Neuroticism expression---records 0.000 across both conditions and all 4,080 runs. LLMs trained through RLHF systematically refuse to generate self-doubt expressions (``I'm probably wrong'', ``I could be mistaken'') regardless of the assigned Neuroticism prior. This constitutes quantitative evidence of an RLHF-imposed behavioral ceiling: the engine cannot improve what the underlying LLM refuses to produce.


\section{Outcome Fidelity Analysis}

Outcome fidelity measures how closely the simulated deliberation outcomes match documented real-world outcomes for each scenario, evaluated across six dimensions: resolution accuracy, stakeholder positioning, dynamics patterns, turning points, world-state changes, and action distribution.

Table~\ref{tab:fidelity} summarizes the outcome fidelity comparison.

\begin{table}[t]
\centering
\caption{Outcome Fidelity: Engine vs.\ Naive}
\label{tab:fidelity}
\begin{tabular}{lcc}
\toprule
Metric & \engine & \naive \\
\midrule
Mean Fidelity             & \textbf{0.302} & 0.161 \\
Welch's $t$-test          & \multicolumn{2}{c}{$t = 3.588$, $p = 0.0004$} \\
Cohen's $d$               & \multicolumn{2}{c}{$+1.135$ (large)} \\
95\% CI ($\Delta$)        & \multicolumn{2}{c}{$[+0.064, +0.218]$} \\
\midrule
Guided fidelity           & \multicolumn{2}{c}{0.270} \\
Exploratory fidelity      & \multicolumn{2}{c}{0.193} \\
\bottomrule
\end{tabular}
\end{table}

The engine achieves 87.6\% higher outcome fidelity than the naive baseline (0.302 vs.\ 0.161), with a large effect size (Cohen's $d = 1.135$) and high statistical significance ($p = 0.0004$). This is the largest performance gap among all measured metrics, suggesting that the action layer and structured world-state tracking contribute substantially to outcome realism beyond what persona fidelity alone can achieve.

Guided simulations (where phase outcomes are specified in the script) produce higher fidelity than exploratory simulations (0.270 vs.\ 0.193), confirming that outcome anchoring provides useful constraints. However, even exploratory simulations under the engine substantially outperform naive simulations (0.193 vs.\ 0.161 for exploratory-only comparison).

Table~\ref{tab:scenario_fidelity} presents the per-scenario fidelity breakdown, ranked by engine performance.

\begin{table}[t]
\centering
\caption{Per-Scenario Outcome Fidelity (Top 5 and Bottom 5)}
\label{tab:scenario_fidelity}
\small
\begin{tabular}{lcc}
\toprule
Scenario & \engine & \naive \\
\midrule
\multicolumn{3}{l}{\textit{Top 5 (Engine)}} \\
Boeing 737 MAX          & 0.581 & 0.372 \\
Fukushima Nuclear       & 0.493 & 0.389 \\
Singapore HDB           & 0.487 & 0.267 \\
Microsoft-Activision    & 0.448 & 0.151 \\
Japan Intern Reform     & 0.414 & 0.425 \\
\midrule
\multicolumn{3}{l}{\textit{Bottom 5 (Engine)}} \\
Starbucks Unionization  & 0.149 & 0.066 \\
Peloton Demand Cliff    & 0.155 & 0.079 \\
FTX Collapse            & 0.161 & 0.072 \\
Theranos Whistleblower  & 0.176 & 0.066 \\
WeWork IPO              & 0.185 & 0.068 \\
\bottomrule
\end{tabular}
\end{table}

The engine achieves universal advantage: it outperforms or matches the naive baseline on 19 of 20 scenarios (the exception being Japan Intern Training Reform, where the naive baseline achieves 0.425 vs.\ the engine's 0.414, a non-significant difference). The scenarios with the highest engine fidelity (Boeing 737 MAX, Fukushima Nuclear) are those with well-documented stakeholder dynamics and clear multi-party negotiation patterns, suggesting that the engine benefits from scenarios where LLM pre-training knowledge provides rich contextual grounding for the action layer.

The scenarios with the lowest fidelity scores (Starbucks Unionization, FTX Collapse, Theranos) are characterized by adversarial dynamics that the engine systematically de-escalates. The outcome analysis reveals a 100\% de-escalation bias: every adversarial scenario is simulated as non-adversarial, with mean tension of 0.000 in expected-high-tension contexts. This reflects the RLHF sycophancy problem manifesting at the outcome level.


\section{Actor Scaling Analysis}

Table~\ref{tab:scaling} presents persona drift and diversity metrics across three actor scales (3, 5, and 10 agents).

\begin{table}[t]
\centering
\caption{Actor Scaling Analysis}
\label{tab:scaling}
\begin{tabular}{lccc}
\toprule
Metric & 3-Actor & 5-Actor & 10-Actor \\
\midrule
\multicolumn{4}{l}{\textit{Engine}} \\
Drift MAE          & 0.171 & 0.165 & 0.168 \\
Envelope Violations & 2.009 & 2.002 & 2.155 \\
Role Diversity      & 0.636 & 0.561 & 0.339 \\
Convergence (mean)  & 0.657 & 0.650 & 0.816 \\
\midrule
\multicolumn{4}{l}{\textit{Naive}} \\
Drift MAE          & 0.183 & 0.176 & 0.176 \\
Envelope Violations & 2.300 & 2.222 & 2.324 \\
\bottomrule
\end{tabular}
\end{table}

\textbf{Persona drift is robust to scaling.} The engine maintains drift between 0.165 and 0.171 across all actor scales, with the maximum increase from 3-actor to 10-actor being only +0.002 (from the best 5-actor result of 0.165 to 0.168 at 10-actor). This demonstrates that the Persona State Controller's per-actor trait tracking scales effectively---each actor's behavioral fidelity is monitored independently, and the addition of more actors does not degrade individual monitoring quality.

\textbf{Role diversity decreases substantially with scale.} Diversity drops from 0.636 (3-actor) to 0.339 (10-actor), a 46.7\% reduction. This is a structural effect: with 10 actors sharing 12 action types, the probability of role overlap increases mechanically. The action family convergence rate confirms this pattern, rising from 0.657 (3-actor) to 0.816 (10-actor). This suggests that the current 12-action vocabulary is adequate for small-group simulations (3--5 actors) but may need expansion for larger-scale deployments.


\section{Phase-Level Analysis}

Table~\ref{tab:phase} presents per-phase drift, convergence, and diversity metrics under the engine condition.

\begin{table}[t]
\centering
\caption{Phase-Level Quality Metrics (Engine Condition)}
\label{tab:phase}
\begin{tabular}{lccc}
\toprule
Phase & Drift & Convergence & Diversity \\
\midrule
OPENING     & 0.171 & 0.741 & 0.357 \\
TENSION     & 0.173 & 0.723 & 0.452 \\
NEGOTIATION & 0.171 & 0.618 & 0.503 \\
CLOSING     & 0.179 & 0.582 & 0.583 \\
\midrule
\multicolumn{4}{l}{\textit{Naive Condition}} \\
OPENING     & 0.179 & --- & --- \\
TENSION     & 0.179 & --- & --- \\
NEGOTIATION & 0.178 & --- & --- \\
CLOSING     & 0.183 & --- & --- \\
\bottomrule
\end{tabular}
\end{table}

Two notable patterns emerge from the phase-level analysis.

\textbf{CLOSING exhibits the highest drift.} Under the engine condition, CLOSING-phase drift (0.179) exceeds all other phases by 0.006--0.008 units. Under the naive condition, the same pattern holds (0.183 vs.\ 0.178--0.179). This confirms the ``summarize and agree'' hypothesis: during the CLOSING phase, actors default to consensus-oriented language patterns that suppress individual personality expression. The engine's phase-specific weight modifiers (Table~\ref{tab:phase_weights}) allocate increased persona and envelope weights during CLOSING ($1.3\times$ and $1.4\times$ respectively) to counteract this tendency, but the drift increase persists, indicating that the RLHF-driven convergence pressure exceeds the engine's corrective capacity.

\textbf{Diversity increases monotonically across phases.} Diversity rises from 0.357 (OPENING) to 0.583 (CLOSING), contrary to the naive expectation that actors would homogenize over time. This reflects the action layer's contribution: as the deliberation progresses from position-stating (OPENING) through conflict (TENSION) to resolution (NEGOTIATION, CLOSING), actors increasingly differentiate their actions---some committing resources, others requesting evidence, others deferring decisions---even as their language converges. The action layer thus provides a structural mechanism for maintaining behavioral diversity that operates independently of linguistic convergence.


\section{Statistical Significance}

Table~\ref{tab:significance} summarizes the statistical tests for primary metrics.

\begin{table}[t]
\centering
\caption{Statistical Significance Summary}
\label{tab:significance}
\begin{tabular}{lccc}
\toprule
Metric & $t$-statistic & $p$-value & Cohen's $d$ \\
\midrule
Persona Drift        & $-10.24$ & $<0.0001$ & $-0.94$ \\
Envelope Violations  & $-4.52$  & $<0.0001$ & $-0.41$ \\
Outcome Fidelity     & $+3.59$  & $0.0004$  & $+1.14$ \\
Fulfillment Rate     & $+3.82$  & $0.0002$  & $+0.35$ \\
Dialogue Coherence   & $+5.17$  & $<0.0001$ & $+0.47$ \\
\bottomrule
\end{tabular}
\end{table}

All primary metrics achieve statistical significance at $p < 0.001$, comfortably exceeding the Bonferroni-corrected threshold of $\alpha' = 0.0005$ for 100 comparisons. The effect sizes range from medium (envelope violations, $d = -0.41$; fulfillment rate, $d = +0.35$) to large (persona drift, $d = -0.94$; outcome fidelity, $d = +1.14$), indicating that the engine's improvements are not merely statistically detectable but practically meaningful.

The non-overlapping 95\% confidence intervals for persona drift provide an intuitive interpretation of significance: the engine's best-case drift (upper bound: 0.1685) is lower than the naive baseline's worst-case drift (lower bound: 0.1773), meaning that even under maximally unfavorable sampling conditions, the engine outperforms the baseline.
